\section{Conclusion}

\label{sec:conclusion}

We have presented the Unified Agent Harness (UAH), a vertically integrated
seven-layer architecture for autonomous AI agents. The central argument of
this paper is both empirical and architectural: the harness, not the model,
is the product.

The empirical case is demonstrated by the competitive analysis: the UAH
achieves 74.3\% on SWE-bench Verified, 70.2\% better absolute performance
than the next best multi-domain competitor, while maintaining competitive
per-task cost. The performance advantage derives not from a superior
underlying model but from cross-layer optimization enabled by vertical
integration.

The architectural case is established by Theorem~\ref{thm:dominance}: any
system that separately optimizes individual layers will yield strictly
suboptimal end-to-end performance compared to joint optimization over all
layers---under the empirically verified condition that layer performances
are non-separable. The seven cross-layer effects we document empirically
confirm this non-separability in production.

The broader implication is organizational as much as technical. As frontier
model capabilities converge (closing the gap between providers that
independently optimize LLMs), the durable competitive advantage in
AI-powered products will come from the infrastructure layer: routing,
tooling, compute, observability, and self-improvement. The AWS analogy
holds: just as cloud infrastructure became the durable business layer above
a commoditizing hardware market, the agent harness will become the durable
business layer above a commoditizing model market.

The UAH defines what a complete agent harness must contain. We release the
architecture specification, the formal layer definitions, and the
open-source components (\texttt{@hanzo/mcp}, \texttt{hanzo-agent},
\texttt{hanzo-operative}) to the community, and invite the field to build
upon, critique, and extend this foundation.

% ============================================================================
% APPENDIX
% ============================================================================
\appendix

